% Pursuing Stacks (cotes 134-2 to 134-8): archive cote 134-8, site page 34 (PDF page 35), typescript page 588.
% typescript. Draft compared: ps7.tex:1594-1658.
% AI-transcribed by Claude Opus 5.5 (claude-opus-5-5, adaptive thinking, effort high), Message Batches API, batch ps4 (msgbatch_012Aoc1NuskATPsyBeKSYK7J), 2026-11-12; not yet checked by eye.
\scanpage{134-8/34}
\mspage{588}

or $Rf_{*}(K^{\bullet})$,  for a general map \xout{f: A B in (Cat) and}
\[ f:\ A \longrightarrow B \]
in (Cat),  and  \xout{L} $F_{\bullet}\in D^{-}(A^{\wedge}_{k})$,  $K^{\bullet}\in D^{+}(A^{\wedge}_{k})$,  I do not see any means of
\margin{L.}bypassing projective resolutions of $\unsure{L}_{\bullet}$ , injective ones of $K^{\bullet}$.  Taking
the more familiar case of $Rf_{*}(K^{\bullet})$,  one natural idea of course would
be to take a projective resolution \xout{of} $L^{A}_{\bullet}$ of $k_{A}$ (i.e. a cointegrator for
A),  and write  tentatively
\[ (1) \qquad Rf_{*}(K^{\bullet}) \overset{?}{\simeq} f_{*}(\underline{Hom}^{\bullet\bullet}_{k}(L^{A}_{\bullet},K^{\bullet}))\ . \]
The (misleading) reflex inducing us to write down this formula, is that
this formula looks as if it were to boil down to the similar (correct)
formula for the maps $A_{/b}\longrightarrow e$ ,  when taking the localizations on B
\[ B_{/b}\longrightarrow B\ . \]
For this intuition to be correct,  \xout{would} it should be true that the
restriction\inserted[typed]{(or ``localization'')} of  $L^{A}_{\bullet}$  to $A_{/b}$  (which is of course a resolution of the
constant sheaf k on $A_{/b}$ ) is indeed a cointegrator on $A_{/b}$ ,  namely
that it's components are still projective.  This however,  we suspect,
will hold true only under very special assumptions  -  as in general,
it is inverse image\xout{s} by \uline{colocalization} (not by localization) which
takes projectives into projectives.  Thus I do'nt expect a relation (1)
to hold,  except under most exacting conditions on f and $K^{\bullet}$,  which I
did'nt try to pin down.  To take an example,  assume A has a final
object hence $k_{A}$ is projective and we may take $L^{A}_{\bullet}=k_{A}$ ,  then (1)
reads ( when $K^{\bullet}=K$ is reduced to degree zero)
\[ Rf_{*}(K)\hand{\simeq \unsure{f}_{*}}(K) \text{ i.e. } R^{i}f_{*}(K)=0 \text{ for } i \hand{>} 0, \]
which need not hold true even if \xout{K is} K is constant ( = $k_{A}$ say).  For
instance,  we may start with an arbitrary object $A_{0}$ of (Cat) and add
a final object\inserted{$e_1$} to get A (intuitively,  it is the \hand{``}cone\hand{''} over $A_{0}$),  which
is mapped into\inserted{the cone over e,} $B = \hand{\Delta}_{1} = (0 \hand{\to} 1)$\hand{,} in the obvious way.  Then for any
sheaf K on A,  with restriction $K_{0}$ to $A_{0}$ ,  we get
\[ R^{i}f_{*}(K)_{0} \text{ (fiber at 0) } = H^{i}(A_{0},K_{0}) \qquad , \]
which needs not be 0 for $i \hand{>} 0$.

There is a big blunder at the end of section 101,  where the
formulae (12),(12'),(13) (p. 377)  (supposed to be ``essentially trivi\xout{al}%
al''),  are false for essentially the same kihd of reason.  The\deleted{se} \deleted{latter}
\deleted{two} formulae r\over{u}{a}n into the typewriter as a matter of course,  as \deleted{they}

\margin{34}
